\documentclass{../latex/utmreport}
\newcommand{\ReportTitle}{Lab 2: Meet the OS You Will Build}
% Shared title-page and PDF metadata for every laboratory report.
\newcommand{\StudentName}{Gafenco Victor}
\newcommand{\StudentGroup}{FAF-241}
\newcommand{\CourseName}{Operating Systems}
\newcommand{\InstructorName}{Patricia Reitman}
\newcommand{\InstructorTitle}{Assistant lecturer}
\newcommand{\ReportCity}{Chișinău}
\newcommand{\ReportYear}{2026}
\newcommand{\CourseRepository}{https://github.com/yoda-dinmd/os-labs}

\selectlanguage{english}
\utmlogo{../latex/assets/UTM_Logo_English.jpeg}
\institution{Technical University of Moldova}
\faculty{Faculty of Computers, Informatics, and Microelectronics}
\department{Department of Software Engineering}
\studyprogramme{Software Engineering Study Programme}
\documenttype{Laboratory report}
\student{\StudentName}{\StudentGroup}
\supervisor{\InstructorName}{\InstructorTitle}
\renewcommand{\UTMCoordinatorLabel}{Instructor}
\cityyear{\ReportCity}{\ReportYear}
\reportsubtitle{\CourseName}
\hypersetup{pdfauthor={\StudentName},pdfsubject={\CourseName},
  pdftitle={\ReportTitle}}
\reporttitle{\ReportTitle}
\addto\captionsenglish{\renewcommand{\figurename}{Figure}\renewcommand{\tablename}{Table}\renewcommand{\bibname}{References}}
\renewcommand{\figurename}{Figure}
\renewcommand{\tablename}{Table}
\renewcommand{\bibname}{References}
\usepackage{float}
\usepackage{placeins}
\setlength{\emergencystretch}{2em}
\newcommand{\code}[1]{\texttt{\detokenize{#1}}}
% Keep screenshots beside their observations, with bounded height and width.
\newcommand{\screenshot}[4][\textwidth]{%
  \begin{figure}[H]
    \centering
    \includegraphics[width=#1,height=0.62\textheight,keepaspectratio]{screenshots/#2}
    \caption{#3}\label{#4}
  \end{figure}}
\newcommand{\reportreferences}[2]{%
  \clearpage\phantomsection\addcontentsline{toc}{chapter}{References}
  \begin{thebibliography}{9}
    \bibitem{handout} Technical University of Moldova, FCIM / FAF.
    \textit{#1}. Operating Systems laboratory handout, academic year 2026--2027.
    Supplied course document: \code{#2}.
  \end{thebibliography}}


\begin{document}
\makeutmreporttitlepage
\tableofcontents
\clearpage

\begin{utmabstract}[Abstract]
This laboratory explores xv6 through compilation, booting, shell commands,
source inspection, and the addition of a user program. The operating system
was run on emulated RISC-V hardware using QEMU from an Ubuntu virtual
machine. Existing commands demonstrated directory listing, file reading,
output, and pipelines. Source inspection connected a user program's read
and write requests with their kernel handlers. A new program named
\code{sleep} was registered in the build and used the \code{pause} system
call declared by the downloaded checkout. Runtime screenshots show a valid
call returning to the shell and usage messages for missing and extra
arguments. The report distinguishes these observed results from unmeasured
properties such as pause duration and exit status.
\utmkeywords{xv6, system calls, user programs, kernel, pipes, emulation}
\end{utmabstract}

\chapter*{Abbreviations}
\addcontentsline{toc}{chapter}{Abbreviations}
\noindent\begin{tabularx}{\textwidth}{lY}
OS & operating system \\
VM & virtual machine \\
PID & process identifier \\
CPU & central processing unit \\
QEMU & the emulator used to run xv6 \\
RISC-V & the instruction-set architecture targeted by this xv6 checkout \\
\end{tabularx}

\utmintroduction
This laboratory follows the course exercise on building and using a small
teaching operating system, xv6~\cite{handout}. It complements the external
inspection of Linux in Lab 1 by exposing both the programs that request OS
services and the kernel code that handles those requests.

The objectives were to boot xv6, use its shell, distinguish userspace from
kernel code, locate system-call handlers, and add a working user program.
Commands were executed in the Ubuntu terminal and inside the xv6 shell.
The figures are the student's captured outputs, and the program listing
matches the source displayed in the VM.

\chapter{Build environment and boot}
The xv6 source tree was obtained in \code{~/xv6-riscv}. The build environment
uses a RISC-V cross-compiler to generate code for the emulated machine.

\section{Source tree and toolchain}
The directory listing contains \code{kernel}, \code{user}, \code{mkfs},
\code{Makefile}, and \code{README}. The tool check located
\code{/usr/bin/riscv64-linux-gnu-gcc}, while the alternative
\code{riscv64-unknown-elf-gcc} check produced no path. QEMU reported version
10.2.1. The header \code{user/user.h} declared \code{int pause(int);}
at line 25, establishing the system-call name required by this checkout.
Figure~\ref{fig:setup} records the checks.
\screenshot{L2_01_setup.png}{Source tree, compiler path, QEMU version, and pause declaration}{fig:setup}

\section{Booting in QEMU}
The command \code{make qemu} launched the kernel. The console printed
\code{xv6 kernel is booting}, messages from additional harts, and
\code{init: starting sh} before displaying the shell prompt.
A hart is a RISC-V hardware thread. Figure~\ref{fig:boot} demonstrates
successful booting. This capture was made after the initial build, so it
shows the launch command rather than a fresh compilation transcript.
\screenshot{L2_02_boot.png}{xv6 boot messages and interactive shell}{fig:boot}

\chapter{Using xv6 as a Unix system}
The shell provides a small set of familiar Unix programs. Their outputs
demonstrate basic file operations and communication between processes.

\section{Included programs and file contents}
Three supplied programs are \code{ls}, which lists directory entries;
\code{cat}, which reads file contents; and \code{echo}, which prints its
arguments. The directory listing also includes \code{grep}, \code{wc},
\code{mkdir}, and \code{sh}. Figure~\ref{fig:programs} shows the listing
and the opening part of the README. The README output is an excerpt,
whereas the separate word-count command operates on the entire file.
\screenshot{L2_03_programs.png}{Included xv6 programs and an excerpt of README}{fig:programs}

\section{Pipelines and shell comparison}
The command \code{echo hello xv6} printed \code{hello xv6}.
The pipeline \code{ls | grep c} produced entries for \code{cat},
\code{echo}, \code{wc}, \code{sync}, and \code{console}.
The command \code{wc README} printed \code{48 336 2441 README}, representing
48 lines, 336 words, and 2441 bytes. Figure~\ref{fig:shell} records these
results.

A pipeline needs process creation and execution, together with interprocess
communication. Separate processes can be started with \code{fork} and
\code{exec}; a kernel pipe transports bytes through file descriptors from
one program's standard output to another's standard input. The xv6 shell
supports familiar execution and pipes, but provides fewer conveniences than
the Linux shell used in Lab 1.
\screenshot[0.75\textwidth]{L2_04_shell.png}{Output, pipeline filtering, and README counts}{fig:shell}

\chapter{Reading userspace and kernel source}
The source tree separates ordinary programs in \code{user/} from privileged
OS code in \code{kernel/}. This distinction makes it possible to trace
a request from a user program to its kernel implementation.

\section{User programs and system calls}
The first 40 lines of \code{user/ls.c} show its includes, filename-formatting
helper, and the opening of a path. The complete \code{user/cat.c} listing
shows how bytes are read into a buffer and written to standard output.
Figures~\ref{fig:ls-source} and~\ref{fig:cat-source} record the inspected code.
\screenshot{L2_05_source_ls.png}{Beginning of the xv6 directory-listing program}{fig:ls-source}
\screenshot{L2_06_source_cat.png}{Complete source of the xv6 cat program}{fig:cat-source}

Table~\ref{tab:calls} identifies the direct system calls visible in
\code{cat.c}. With no file arguments, the program reads standard input,
descriptor 0. Its output goes to descriptor 1.
\begin{table}[H]
\caption{System calls used by the cat program}\label{tab:calls}
\centering
\begin{tabularx}{\textwidth}{|l|Y|}\hline
\textbf{Call} & \textbf{Request to the kernel}\\\hline
\code{open} & Open an input pathname and obtain a descriptor.\\\hline
\code{read} & Transfer input bytes into the userspace buffer.\\\hline
\code{write} & Transfer bytes from the buffer to standard output.\\\hline
\code{close} & Release the opened descriptor.\\\hline
\code{exit} & Terminate the process with a success or error status.\\\hline
\end{tabularx}
\end{table}
The function \code{fprintf} formats output in userspace; it is not itself
a system call, although its output ultimately uses \code{write}.
Figure~\ref{fig:declarations} shows the system-call declarations followed by
userspace library declarations in \code{user/user.h}.
\screenshot{L2_07_syscall_declarations.png}{System-call interface and userspace library declarations}{fig:declarations}

\section{Kernel handlers and privilege separation}
In the inspected checkout, \code{sys_read(void)} is at line 69 of
\code{kernel/sysfile.c}; its return type appears on line 68.
The handler retrieves arguments, checks the file descriptor, and calls
\code{fileread}. The handler \code{sys_write(void)} is at line 83 and calls
\code{filewrite}. Figure~\ref{fig:handlers} records both implementations.
These line numbers refer to the captured source version.

The code in \code{kernel/} runs with privilege to manage resources and
implement OS services, while the code in \code{user/} runs as ordinary
programs and requests those services through system calls.
\screenshot{L2_08_syscall_kernel.png}{Kernel implementations of read and write}{fig:handlers}

\chapter{Implementing and testing sleep}
The new program pauses for a requested number of clock ticks and returns to
the shell. Its executable name is \code{sleep}, while its kernel request
uses the \code{pause} interface declared by this checkout.

\section{Program source and version adaptation}
The handout's example calls \code{sleep(atoi(argv[1]))}. The inspected
\code{user/user.h} declares \code{pause} instead, so the implementation
calls \code{pause(atoi(argv[1]))}. This adapts the system-call name while
preserving the executable and usage syntax required by the exercise.

The program requires exactly one command-line argument. Otherwise it prints
\code{usage: sleep <ticks>} to standard error and calls \code{exit(1)}.
For a valid argument count, \code{atoi} converts the argument to an integer,
\code{pause} requests the delay, and \code{exit(0)} terminates the program.
The source is shown below and in Figure~\ref{fig:sleep-source}.
\lstinputlisting[language=C]{user/sleep.c}
\screenshot{L2_09_sleep_code.png}{Creating and displaying the sleep source in the VM}{fig:sleep-source}

\section{Build registration}
The executable was added to the \code{UPROGS} list as follows:
\begin{lstlisting}
  $U/_sync\
  $U/_sleep\

fs.img: mkfs/mkfs README $(UPROGS)
\end{lstlisting}
The actual Makefile entries use tab indentation. A blank line is preserved
before \code{fs.img:} because the final list entry ends with a continuation
backslash. Figure~\ref{fig:makefile} records an earlier edit in which this
separator had been removed; that capture is not a representation of the
correct final layout. The repository includes \code{Makefile.patch} with
the intended registration and separator.
\screenshot{L2_10_makefile.png}{Earlier Makefile edit before restoring the blank separator}{fig:makefile}

\section{Observed runtime results}
The runtime sequence printed \code{before}, ran \code{sleep 10}, and then
printed \code{after}. The valid call returned to the shell without an error
message. Running \code{sleep} without arguments printed the usage message.
Figure~\ref{fig:sleep-test} shows these results and an earlier failed
extra-argument attempt.
\screenshot[0.60\textwidth]{L2_11_sleep_test.png}{Valid call, missing-argument check, and earlier failed attempt}{fig:sleep-test}

The repeated command \code{sleep 10 20} printed
\code{usage: sleep <ticks>}, as shown in Figure~\ref{fig:extra-args}.
Both invalid argument counts therefore produced the expected message.
The earlier \code{exec sleep failed} result was superseded by this successful
repeat; its cause was not established by the captures.
\screenshot[0.55\textwidth]{fixed2.png}{Successful extra-argument usage check}{fig:extra-args}

The screenshots show execution and argument-count handling, but do not
measure the pause duration or expose exit statuses. No independent rebuild
was performed from this report repository because it contains the user
program and registration patch rather than the full xv6 source tree.

\utmconclusions
xv6 booted successfully in QEMU and provided working file-reading, output,
and pipeline commands. Inspecting both \code{user/cat.c} and
\code{kernel/sysfile.c} connected the user-level read and write requests
with their privileged kernel handlers.

The added \code{sleep} program uses the checkout's \code{pause} system
call and is accompanied by a build-registration patch. The captured runtime
results show a valid call returning and the expected usage messages for
missing and extra arguments. The repository containing the reports,
screenshots, and program is \expandafter\url\expandafter{\CourseRepository}.

\end{document}
